\documentclass[11pt]{article}

\usepackage[margin=2.2cm]{geometry}
\usepackage{xcolor}
\usepackage{forest}
\usepackage{tikz}
\usetikzlibrary{arrows.meta,positioning,decorations.pathreplacing,calc}
\usepackage{amsmath,amssymb,mathtools}
\usepackage{booktabs}
\usepackage{array}
\usepackage{enumitem}
\usepackage{titlesec}
\usepackage{float}
\usepackage[colorlinks=true,linkcolor=teal!70!black,citecolor=teal,urlcolor=teal]{hyperref}
\usepackage{parskip}
\usepackage{cleveref}

\setlength{\parskip}{0.5\baselineskip}
\titlespacing*{\section}{0pt}{1.5em}{0.6em}
\titlespacing*{\subsection}{0pt}{1.2em}{0.4em}

% ---- automatic "Part N --- Title" section numbering, so \section + \label
% gives both the on-page heading and a working \ref/\cref everywhere else ----
\renewcommand{\thesection}{Part \arabic{section}}
\titleformat{\section}{\normalfont\Large\bfseries}{\thesection}{0.6em}{---\ }

% ---- colors reused across every figure, so the same idea always wears the same color ----
\definecolor{lossless}{HTML}{1f7a6c}   % teal family
\definecolor{lossy}{HTML}{c2622d}      % orange family
\definecolor{prefixcol}{HTML}{6a4c93}  % violet
\definecolor{nonprefixcol}{HTML}{1f6fb2}% blue
\definecolor{dictcol}{HTML}{8a5a2b}    % brown
\definecolor{bestcol}{HTML}{2e8b3d}    % green, "best in class" marker
\definecolor{warncol}{HTML}{b3261e}    % red, "error / impossible" marker
\definecolor{transformcol}{HTML}{2f7d8f}% steel teal, transform-based family
\definecolor{contextcol}{HTML}{a13d63} % magenta, context-modeling family
\definecolor{mmcol}{HTML}{4f6d2f}      % olive, lossless multimedia family
\definecolor{groupcol}{HTML}{1f6fb2}   % blue, used for the group axioms box

\newcommand{\figdrawnote}[1]{\par\noindent\textbf{\textcolor{gray!60!black}{[Draw on the board here:}} \textit{#1}\textbf{\textcolor{gray!60!black}{]}}\par}

% ---- a colored callout box for direct-to-camera, pause-and-think moments ----
\newcommand{\engage}[3]{%
\begin{center}
\begin{tikzpicture}
\node[draw=#1, thick, rounded corners, fill=#1!6, align=left,
      font=\sffamily\small, text width=0.92\linewidth, inner sep=9pt]
      {\textbf{\textcolor{#1}{#2}}\; #3};
\end{tikzpicture}
\end{center}
}

% ---- a light gray box for verbatim "say this to camera" narration ----
\newcommand{\say}[1]{%
\begin{center}
\begin{tikzpicture}
\node[draw=gray!50, thick, rounded corners, fill=gray!5, align=left,
      font=\itshape, text width=0.92\linewidth, inner sep=9pt]
      {#1};
\end{tikzpicture}
\end{center}
}

\title{\Large\bfseries Tiling a Lattice by Transforming its Fundamental Region\\[2pt]
\normalsize\normalfont\itshape A visual definition of lattice}
\author{Dr. Faisal Aslam}
\date{}

\begin{document}
\maketitle
\vspace{-2em}

In the previous video, we looked at two definitions of a lattice. In this video, we will look at another definition that can be used to develop a visual understanding of a lattice. Furthermore, this definition is also helpful in proving many theorems related to lattices. Without further ado, let's define a lattice using the concept of a fundamental region.

A lattice $L \subset \mathbb{R}^n$ has a fundamental domain (also called a fundamental region) such that by translating this fundamental domain by lattice points, we can tile the whole of $\mathbb{R}^n$. You can think of this fundamental region as a basic tile that can be used to tile the whole of $\mathbb{R}^n$ using lattice points. Furthermore, this tile has a fixed volume, which is referred to as the volume of the lattice. So in this video, we will first provide mathematical definitions of a) fundamental region, b) volume of the lattice, and c) how to shift that region to tile $\mathbb{R}^n$. You might not understand these definitions readily; therefore, we will then go through several examples so that these three concepts are crystal clear. So let's start with the mathematical definition.

\paragraph{Definition 1: Fundamental Region}

For a lattice generated by basis vectors $\mathbf{b}_1, \mathbf{b}_2, \dots, \mathbf{b}_n$, the fundamental region is the \textbf{fundamental parallelepiped}:
$$
\mathcal{P} = \left\{ \sum_{i=1}^{n} \alpha_i \mathbf{b}_i : 0 \le \alpha_i < 1 \right\}.
$$

For the rest of this video, we will refer to the fundamental region as the fundamental parallelepiped.

\paragraph{Definition 2: Volume of the Lattice}

The volume of a lattice $L$ with basis vectors $\mathbf{b}_1, \mathbf{b}_2, \dots, \mathbf{b}_n$ is
$$
\text{Vol}(L) = \left| \det(\mathbf{b}_1, \mathbf{b}_2, \dots, \mathbf{b}_n) \right|.
$$

\paragraph{Definition 3: Tiling $\mathbb{R}^n$}

Finally, the last definition is about shifting our basic tile in such a manner that it covers the whole of $\mathbb{R}^n$. We can simply shift our fundamental parallelepiped (i.e., our basic tile) by adding a lattice point to it. That shifting is simply $\mathcal{P} + l$ where $l \in L$.

If you have understood the above concepts, then I congratulate your intelligence. However, my aim is that everyone can understand them and everything is crystal clear, so I will now explain these concepts with some examples.

\end{document}
